\chapter{Algebraic Spaces}

We have seen in ... that the representable presheaf of a scheme is indeeed a sheaf in $\SchOver{Et}$. Via the Yoneda embedding we can identify a scheme $U$ with its correspoding sheaf and denote it also by $U$.

\begin{definition}[Representable Morphism]
	(\cite[][(Definition 4.1.1), p.~105]{alper2026stacks})
	A morphism $X\to Y$ of sheaves over $\Sch$ is called representable by schemes, if for every scheme $T$ and every morphism $T\to Y$, the fibre product $X\times_Y T$ is a scheme.
\end{definition}

\begin{definition}[\Etale\ morphisms of sheaves]
	Given a morphism $f:X\to Y$ of sheaves over the big \etale\ site. We call $f$ \etale, if $f$ is representable and for every scheme $T\to Y$, the projection $X\times_YT\to T$ is \etale.
\end{definition}

\begin{remark}
	Note, that via Yoneda the morohism $X\times_Y T \to T$ corresponds to a morphism of scheme, since both are representable by assumption.
\end{remark}

\begin{definition}[Algebraic Space]
	(\cite[][(Definition 4.1.2), p.~105]{alper2026stacks})
	An algebraic space $X$ is a sheaf $X:\SchOver{Et} \to \Set$, such that there exists a scheme $U$ and an \etale\ surjection $U\to X$ representable by schemes.
\end{definition}

\begin{remark}
	In analogy with schemes, which are Zariski-locally given by affine schemes, an algebraic space is a sheaf admitting an \etale\ cover by schemes, and hence, after refinement (\ref{prop:affine-etale-refinement}), by affine schemes. Thus an algebraic space is a sheaf that is \etale-locally a scheme.
	
	Indeed, let $U \to X$ be an \etale\ surjection, where $U$ is a scheme, and let $\{U_i\}$ be an open affine cover of $U$. Then the induced morphism
	\[
	\coprod_i U_i \to X
	\]
	is again an \etale\ surjection. It follows that every algebraic space admits an \etale\ cover by affine schemes.
\end{remark}



\begin{theorem}[Representability of the Diagonal]
	(\cite[][(Theorem 4.2.1), p.~118]{alper2026stacks})
	The diagonal $\Delta_X: X\to X\times X$ of an algebraic space $X$ is represantable by schemes.
\end{theorem}

\begin{definition}[\Etale\ equivalence relation]
	(\cite[add. text]{olsson2016spaces})
	An \etale\ equivalence relation on an $S$-scheme $X$ is a monomorphism 
	\[
		R\hookrightarrow X\times_S X,
	\]
	such that for every $S$-scheme $T$, the $T$-points $R(T) \subset X(T)\times X(T)$ is an equivalence relation on $X(T)$ and the projections $s,t: R\to X$ are \etale.
\end{definition}

